\documentclass[11pt]{article}

\usepackage[margin=1in]{geometry}
\usepackage{amsmath,amssymb}
\usepackage{booktabs}
\usepackage{array}
\usepackage{longtable}
\usepackage{graphicx}

\usepackage{hyperref}
\hypersetup{colorlinks=true, linkcolor=blue, citecolor=blue, urlcolor=blue}
\usepackage[T1]{fontenc}
\usepackage{times}
% Use Latin Modern for typewriter text: the URW Courier (pcr) metrics that
% `times` selects for \ttdefault are not installed in this TeX Live scheme,
% while Latin Modern (lmodern) is always available.
\renewcommand{\ttdefault}{lmtt}

\title{\bfseries BEHRT for Longitudinal Clinical Code Prediction:\\
A Compact Transformer for Next-Visit Diagnosis Forecasting\\
in Multimorbid ICU Patients}

\author{
Fabrisio Braulio Ponte Vela\\
\textit{Frost Institute for Data Science and Computing (IDSC)}\\
\textit{University of Miami, Coral Gables, FL, USA}\\
\texttt{fabrisio.ponte@gmail.com}
}

\date{}

\begin{document}

\maketitle

\begin{abstract}
Longitudinal prediction of which clinical codes will appear at a patient's next
visit is a central problem in clinical informatics, yet most electronic health
record (EHR) models trained on large, general-population cohorts do not
characterize \emph{what kind of task} they are actually solving. We adapt
BEHRT, a BERT-style transformer for EHR sequences, to a cohort of 24{,}729
critically ill patients (179{,}713 training/evaluation samples) drawn from a
MIMIC-IV-derived, ICU population, with diagnoses grouped into 463 Clinical
Classifications Software Refined (CCSR) categories. Our reference model uses
6 transformer layers, 12 attention heads, and a 288-dimensional hidden space,
trained with masked language model (MLM) pre-training followed by multi-label
next-visit fine-tuning with positive-class weighting for rare codes. The model
achieves sample-wise ROC-AUC of 0.901 and Average Precision Score (APS) of
0.274 on held-out test data. Beyond headline metrics, we present eleven
exploratory data analyses and two gradient/attention-based interpretability
studies that characterize \emph{what} the model is actually predicting: 70\%
of next-visit target codes are recurrences of a patient's own diagnosis
history rather than new onsets; positive-class weighting (capped at 30)
improves APS by a verified +7.6\% (0.2543 $\to$ 0.2735) in a matched,
eval-schema-consistent comparison; performance varies sharply by clinical code
type, from F1 = 0.386 on administrative codes to F1 = 0.046 on injury codes;
and longer diagnosis histories improve performance even when truncated to the
model's 64-token context window. Integrated Gradients and attention-rollout
analyses corroborate that a meaningful fraction (53.3\%) of the model's false
positives are driven by genuine cross-category comorbidity signal rather than
superficial code-similarity confusion. Critically, we also compare the model
against two non-learned baselines and two learned, non-sequential baselines
(logistic regression and XGBoost) built on flat per-patient history features:
both learned baselines match or exceed BEHRT on aggregate ranking metrics,
while BEHRT's evidenced advantage is concentrated specifically in
new-diagnosis detection and in producing usable, well-calibrated point
predictions from a single threshold. We argue that this body of evidence
supports a narrow but defensible claim: BEHRT-style transformers can learn
temporal patterns of clinical code recurrence in multimorbid ICU patients, and
should be framed as care-coordination support tools rather than general
disease-onset predictors.
\end{abstract}

\section{Introduction}

Precision healthcare aims to intervene earlier and more precisely by
anticipating a patient's clinical trajectory before it fully unfolds. Deep
learning models trained directly on longitudinal electronic health records
(EHR) have shown that self-attention architectures can learn rich
representations of disease co-occurrence and temporal progression without
hand-crafted features~\cite{li2020behrt,vaswani2017attention}. BEHRT
in particular demonstrated that treating a patient's visit history as a
sequence of diagnosis ``tokens'' -- analogous to words in a sentence, with
visits playing the role of sentence boundaries -- allows a BERT-style
transformer~\cite{devlin2019bert} to be pre-trained with a masked language
model (MLM) objective and then fine-tuned for downstream disease prediction
tasks.

Most published BEHRT-style studies, however, are trained on very large,
general-population cohorts (BEHRT itself used 1.6 million primary-care
patients) and report only aggregate ranking metrics (ROC-AUC, APS) without
characterizing the composition of the prediction target itself. This leaves
an important question unanswered for smaller, higher-acuity cohorts: when a
model reports a high AUC, is it learning to anticipate genuinely new disease
onset, or is it primarily re-identifying chronic conditions a patient already
has? Is performance driven uniformly across disease categories, or is it
concentrated in a handful of common, high-frequency codes while symptom,
administrative, and injury codes remain effectively unpredictable?

In this work we adapt BEHRT to a substantially smaller, ICU-derived cohort
(24{,}729 patients, 179{,}713 samples) using Clinical Classifications Software
Refined (CCSR) diagnosis groupings, and rather than stopping at a single
headline AUC/APS pair, we deliberately interrogate the prediction task itself.
Concretely, our contributions are:

\begin{enumerate}
  \item A BEHRT-style architecture (6 layers, 12 heads, 288-dimensional
    hidden space, 474-entry CCSR vocabulary, 64-token context) trained with
    two-phase MLM pre-training and multi-label next-visit fine-tuning, using
    positive-class weighting to address a 52{,}559:1 label imbalance ratio.
  \item An eleven-step exploratory data analysis pipeline that characterizes
    the dataset's recurrence-vs-onset composition (70\%/30\%), its clinical
    code-type heterogeneity (85.7\% disease, 7.1\% symptom, 4.3\%
    administrative, 2.6\% injury, 0.3\% pregnancy/perinatal), its extreme
    class imbalance, and how performance varies by label support tier and
    diagnosis-history length.
  \item A rigorously verified quantification of the effect of positive-class
    weighting on rare-code detection, including the discovery and correction
    of an eval-schema inconsistency across earlier training runs that would
    otherwise have produced a misleading comparison.
  \item Two interpretability analyses -- Integrated
    Gradients~\cite{sundararajan2017axiomatic} attribution and attention
    rollout~\cite{abnar2020quantifying} -- that test whether the model's false
    positives and attention patterns reflect genuine cross-category
    comorbidity signal, and whether raw last-layer attention under-represents
    long-range dependencies relative to the full attention-flow graph.
  \item An honest, evidence-grounded discussion of what this model can and
    cannot claim, intended to guide the framing of BEHRT-style models trained
    on smaller, higher-acuity cohorts.
\end{enumerate}

\section{Related Work}

BERT~\cite{devlin2019bert} introduced bidirectional masked-language-model
pre-training for natural language and was subsequently adapted to clinical
text by BioBERT~\cite{lee2020biobert} and
ClinicalBERT~\cite{alsentzer2019clinicalbert}. BEHRT~\cite{li2020behrt} applied
the same pre-training paradigm directly to sequences of structured diagnosis
codes rather than free text, using disease, age, segment (visit-separator),
and position embeddings, and demonstrated an 8--10.8 percentage-point absolute
improvement in APS over RETAIN~\cite{choi2016retain} and
Deepr~\cite{nguyen2017deepr} on a 1.6-million-patient primary-care cohort.
Med-BERT~\cite{rasmy2021medbert} extended this line of work with visit-level
representations evaluated on large US insurance-claims data. Our work targets
a materially different regime: a two-orders-of-magnitude smaller, ICU-derived
cohort with 463 active CCSR diagnosis categories, where label scarcity and
class imbalance -- not scale -- are the dominant modeling challenges.

Attention-based interpretability of transformer models has been shown to be
descriptive rather than strictly causal~\cite{jain2019attention}; we therefore
treat raw attention analysis as exploratory and corroborate it with two
complementary techniques: Integrated
Gradients~\cite{sundararajan2017axiomatic}, a gradient-based attribution
method with an axiomatic completeness guarantee, and attention
rollout~\cite{abnar2020quantifying}, which propagates attention through
residual connections across all layers to approximate the full
information-flow graph rather than only the last layer's raw attention.

\section{Materials and Methods}

\subsection{Data Source and Cohort}

The underlying EHR data is derived from a MIMIC-IV-style ICU
database~\cite{johnson2023mimiciv}. After preprocessing (Section~3.2), the
study cohort comprises 24{,}729 patients, each with a minimum of 5 recorded
visits, yielding 179{,}713 labeled training/evaluation samples (median
7 diagnosis-visit ``snapshots'' per patient). The cohort is skewed toward
older, multimorbid adults: median patient age is 62 years, median diagnosis
history length is 44 codes, 88\% of patients are 40 years or older, 46\% are
65 or older, and 0\% are under 18. Each sample's prediction target is a
multi-hot vector over 463 active CCSR diagnosis categories (474-entry
vocabulary including special tokens), representing all diagnoses that appear
at the patient's next visit; the median target size is 9 codes.

\subsection{Preprocessing Pipeline}

Raw ICD-9/ICD-10 codes ($\sim$89{,}000 unique codes) are mapped to CCSR
categories using the CMS/HCUP crosswalk tables~\cite{hcup2021ccsr}. Codes
without a valid mapping are dropped, dot notation is normalized (e.g.
\texttt{I10.9}~$\to$~\texttt{I109}), and visit-level deduplication retains only
unique CCSR codes per visit. A generic catch-all code and placeholder patterns
(matching \texttt{XXX*}, \texttt{000}, \texttt{999}, \texttt{UNK}) are removed
as noise; an independent cleaning-impact audit (Section~5) confirmed this
strategy behaves consistently across the train/validation/test splits.
Patients with fewer than 5 visits, or with zero remaining labels after
cleaning, are excluded. Patient age at each visit is capped at 110 years for
privacy compliance and discretized into an age-embedding vocabulary of 1{,}322
bins. All patient visits are ordered chronologically, and sequences are
truncated to the most recent 64 tokens (the model's maximum position-embedding
length) when a patient's full history exceeds that length.

\subsection{BEHRT Architecture}

Each input token's representation is the sum of four learned embeddings --
diagnosis-code, segment (visit-separator, A/B), age-bin, and sinusoidal
position -- following the original BEHRT
formulation~\cite{li2020behrt,devlin2019bert}. The combined embedding is passed
through a stack of standard post-layer-norm transformer encoder blocks (each
consisting of multi-head self-attention, a residual connection with layer
normalization, a GeLU-activated feed-forward network, and a second residual
connection with layer normalization), followed by a pooling layer over the
\texttt{[CLS]} position and a linear classification head with 463 output
units (sigmoid activation, multi-label). Table~\ref{tab:arch} summarizes the
exact configuration of our reference model.

\begin{table}[h]
\centering
\caption{Reference model configuration
(\texttt{clean\_run\_20260910\_124255}).}
\label{tab:arch}
\begin{tabular}{ll}
\toprule
Parameter & Value \\
\midrule
Transformer layers & 6 \\
Attention heads & 12 \\
Hidden size & 288 \\
Intermediate (feed-forward) size & 256 \\
Max sequence length (context window) & 64 tokens \\
CCSR vocabulary size & 474 (463 active in training data) \\
Age-embedding vocabulary size & 1{,}322 \\
Segment-embedding vocabulary size & 2 \\
Hidden / attention dropout & 0.1 \\
Hidden activation & GeLU \\
Batch size & 64 \\
Optimizer & Adam, gradient-norm clipping at 1.0 \\
Random seed & 42 \\
\bottomrule
\end{tabular}
\end{table}

\subsection{MLM Pre-training}

Following the standard BERT/BEHRT masking recipe, 15\% of diagnosis tokens in
each sequence are selected; of these, 80\% are replaced with a
\texttt{[MASK]} token, 10\% with a random diagnosis code, and 10\% are left
unchanged. The model is trained to predict the original code at each masked
position from bidirectional context, forcing it to learn a joint
representation of diagnosis co-occurrence, visit structure, and age-linked
temporal progression without any downstream supervision signal.

\subsection{Next-Visit Multi-Label Fine-Tuning}

The downstream task is framed as multi-label binary classification: given a
patient's diagnosis history up to and including visit $t$, predict the
multi-hot vector of CCSR codes present at visit $t+1$. The classification head
is trained with a multi-label soft-margin (binary cross-entropy) loss, with an
optional per-class positive weight capped at 30 (\texttt{max\_pos\_weight =
30}) to counteract the severe label imbalance described in Section~5.1. The
final decision threshold is not fixed at the conventional 0.5; instead it is
tuned on a held-out validation split by sweeping thresholds from 0.10 to 0.90
in steps of 0.05 and selecting the value that maximizes micro-averaged F1
(Section~4). A special \texttt{UNK} label is excluded from all reported
metrics, since it does not correspond to a clinically meaningful diagnosis
category.

\section{Training and Evaluation Protocol}

The model is trained for 3 epochs with the configuration in
Table~\ref{tab:arch}, using a fixed seed (42) and the full training split
(no sample-limiting/smoke-test subsampling). Evaluation uses two
complementary, threshold-free ranking metrics computed per sample and
macro-averaged across the test set: sample-wise ROC-AUC (probability that a
true positive code is ranked above a true negative code, for a given patient)
and sample-wise Average Precision Score (area under the precision--recall
curve, which is more sensitive to class imbalance than AUC). For
threshold-dependent metrics (precision, recall, F1), we report results at the
validation-tuned decision threshold (0.60 for the reference model) unless
otherwise noted.

A methodological note that materially affects Section~5.3: over the course of
this project, the evaluation pipeline itself changed twice -- \texttt{UNK}
exclusion and validation-based threshold tuning were introduced partway
through the project's history. Earlier training runs that predate these
changes report inconsistent, sometimes duplicate-looking APS/AUC values for
otherwise-identical configurations. All comparisons in this paper that involve
multiple runs (Section~5.3) are restricted to runs sharing the \emph{current}
evaluation schema; runs under the legacy schema are explicitly excluded and
flagged as such, rather than silently averaged in.

\section{Results}

\subsection{Dataset Characterization}

Three properties of the dataset materially shape how its predictive
performance should be interpreted.

\textbf{Recurrence, not onset.} Of all next-visit target codes, 70\% are
codes already present somewhere in the patient's own diagnosis history (i.e.
recurrences of a known chronic condition), and only 30\% are genuinely new
diagnoses. This means the task is closer to ``which of this patient's known
conditions will be documented again at the next visit'' than to general
disease-onset prediction.

\textbf{Code-type heterogeneity.} Not all CCSR target codes represent
diseases. Based on a category-to-code-type mapping applied to the current
reference run's label vocabulary, target occurrences break down as: 85.7\%
disease diagnoses, 7.1\% symptoms (e.g. abnormal labs, pain), 4.3\%
administrative/encounter codes, 2.6\% injury/trauma codes, and 0.3\%
pregnancy/perinatal-related codes (14.3\% non-disease codes in total). Injury
and symptom codes are, by nature, driven by acute and often unpredictable
events rather than temporal progression, which materially caps achievable
performance in those categories (Section~5.4).

\textbf{Extreme class imbalance.} The single most common target code
(\texttt{CCSR\_END010}, an endocrine/metabolic category consistent with
diabetes) appears in 52{,}559 of 1{,}360{,}288 total training label-occurrences
(30.7\% train-support share among positives at the label level), while the
rarest labels have as few as 1--9 training examples. The resulting positive-
to-negative imbalance ratio is approximately 52{,}559:1 at the extremes, and
only about 2\% of all (sample, label) pairs are positive overall.

\subsection{Overall Next-Visit Prediction Performance}

At the validation-tuned decision threshold of 0.60, the reference model
achieves a sample-wise ROC-AUC of \textbf{0.9006} and a sample-wise Average
Precision Score of \textbf{0.2735} on the held-out test split, with a
micro-averaged F1 of 0.238. The gap between a high AUC and a comparatively
modest APS is expected and informative for a 463-class, severely imbalanced
multi-label problem: AUC measures ranking quality in a class-balanced sense,
while APS is directly penalized by the underlying label imbalance and better
reflects the full precision--recall trade-off a clinician would experience
when reviewing ranked predictions.

To contextualize these numbers, we compared the reference model against four
non-neural baselines evaluated under the identical schema (same label
vocabulary, same UNK exclusion, same threshold-tuning procedure): two
trivial, non-learned baselines (\texttt{repeat\_last\_visit}, which simply
copies forward the codes from the patient's most recent visit, and
\texttt{popularity}, which scores every patient by fixed training-set code
frequencies with no patient-specific information at all) and two learned,
non-sequential baselines built on flat, hand-engineered per-patient history
features (one-vs-rest logistic regression and gradient-boosted trees via
XGBoost, both using the same 942-dimensional feature vector: multi-hot
code-in-history and code-in-last-visit indicators, patient age, and visit
count). Table~\ref{tab:baselines} reports the comparison.

\begin{table}[h]
\centering
\caption{Reference BEHRT vs.\ non-neural baselines (test set, $n=18{,}237$).}
\label{tab:baselines}
\begin{tabular}{lcccc}
\toprule
Model & APS & ROC-AUC & Tuned threshold & Micro-F1 (tuned) \\
\midrule
XGBoost & \textbf{0.4894} & \textbf{0.9307} & 0.90 & 0.147 \\
Logistic regression & 0.2915 & 0.9162 & 0.80 & \textbf{0.297} \\
BEHRT (reference) & 0.2735 & 0.9006 & 0.60 & 0.238 \\
Popularity & 0.2544 & 0.8914 & 0.15 & -- \\
Repeat last visit & 0.2372 & 0.7184 & 0.10 & -- \\
\bottomrule
\end{tabular}
\end{table}

This comparison substantially changes how the reference model's performance
should be interpreted. On both ranking metrics (APS, ROC-AUC), BEHRT is
outperformed by both learned flat-feature baselines: logistic regression
exceeds it by +6.6\% relative APS, and XGBoost by +78.9\% relative APS.
However, ranking quality and point-prediction usability diverge sharply for
XGBoost: despite having the best APS/AUC of any model tested, its micro-F1 at
its own validation-tuned threshold (0.147) is the worst of the three learned
models -- a calibration artifact of using a single dataset-wide
\texttt{scale\_pos\_weight} across 470 labels whose base rates differ by more
than three orders of magnitude, rather than evidence of weaker underlying
signal.

To confirm these ranking-metric gaps are not an artifact of sampling noise in
a single fixed test split, we computed 500-resample bootstrap 95\% confidence
intervals for APS and ROC-AUC for every model, together with paired
significance tests of each baseline's delta against BEHRT (shared resample
indices across models). Table~\ref{tab:bootstrap} reports the result: all
eight paired deltas (APS and AUC, four baselines) are statistically
significant at $p<0.001$, confirming that both the XGBoost/logistic-regression
advantage on ranking metrics and BEHRT's advantage over the two non-learned
baselines are genuine rather than noise.

\begin{table}[h]
\centering
\caption{Bootstrap 95\% confidence intervals (500 resamples, test set) and
significance of each model's paired delta vs.\ BEHRT.}
\label{tab:bootstrap}
\begin{tabular}{lccc}
\toprule
Model & APS (95\% CI) & ROC-AUC (95\% CI) & $p$ vs.\ BEHRT (APS, AUC) \\
\midrule
XGBoost & 0.4894 (0.4863, 0.4924) & 0.9307 (0.9297, 0.9318) & $<0.001$, $<0.001$ \\
Logistic regression & 0.2915 (0.2890, 0.2939) & 0.9162 (0.9153, 0.9171) & $<0.001$, $<0.001$ \\
BEHRT (reference) & 0.2735 (0.2708, 0.2762) & 0.9006 (0.8997, 0.9016) & -- \\
Popularity & 0.2544 (0.2521, 0.2569) & 0.8914 (0.8902, 0.8924) & $<0.001$, $<0.001$ \\
Repeat last visit & 0.2372 (0.2346, 0.2398) & 0.7184 (0.7165, 0.7202) & $<0.001$, $<0.001$ \\
\bottomrule
\end{tabular}
\end{table}

Table~\ref{tab:baselines-newrecur} decomposes each learned model's
performance by whether the target is a recurrence or a new diagnosis.

\begin{table}[h]
\centering
\caption{New-vs-recurring diagnosis performance across learned models.}
\label{tab:baselines-newrecur}
\begin{tabular}{lcccc}
\toprule
& \multicolumn{2}{c}{Recurring} & \multicolumn{2}{c}{New diagnosis} \\
Model & F1 & ROC-AUC & F1 & ROC-AUC \\
\midrule
XGBoost & 0.440 & 0.760 & 0.050 & 0.855 \\
Logistic regression & \textbf{0.453} & 0.676 & 0.025 & 0.838 \\
BEHRT (reference) & 0.403 & 0.681 & \textbf{0.061} & \textbf{0.865} \\
\bottomrule
\end{tabular}
\end{table}

This decomposition is the key to reconciling the aggregate result with a
defensible role for BEHRT. Most of the aggregate metrics are driven by the
majority (70\%) recurring-code subset, where a simple linear model over
which codes already appear in a patient's history captures most of the
available signal (logistic regression's recurring F1 of 0.453 is the
highest of any model tested). On the minority, harder, and arguably more
clinically useful new-diagnosis subset, BEHRT retains the highest F1 (0.061,
roughly 2.4$\times$ logistic regression's) and the highest ranking AUC
(0.865), though XGBoost's new-diagnosis AUC (0.855) is close behind,
indicating a strong flat-feature ranking signal exists there too. We
therefore revise our central claim: BEHRT's evidenced advantage over
comparably simple baselines is concentrated in new-diagnosis detection and in
producing a well-calibrated single operating threshold, not in aggregate
ranking metrics, which a non-sequential model can match or exceed on this
task.

A natural objection to XGBoost's poor point-prediction F1 is that it is
merely a threshold-calibration artifact of using one dataset-wide
\texttt{scale\_pos\_weight}, and that per-label threshold tuning would close
most of the gap to BEHRT. We tested this directly: refitting the identical
XGBoost model and replacing its single global decision threshold (tuned on
validation micro-F1) with 470 independently F1-maximized per-label
thresholds (46 of which had zero validation positives and fell back to the
global threshold) recovers only +0.0079 absolute test micro-F1
(0.1467~$\to$~0.1546) -- a small fraction of the 0.091 gap to BEHRT's 0.238.
This rejects a purely threshold-based explanation and points instead to the
training-time \texttt{scale\_pos\_weight} itself -- a single value averaged
across labels whose base rates differ by more than three orders of
magnitude -- as the more likely source of XGBoost's poorly calibrated
per-label probabilities, rather than the choice of decision threshold alone.

\subsection{Effect of Positive-Class Weighting}

To quantify the effect of positive-class weighting, we compared the reference
model (\texttt{use\_pos\_weight=True}, \texttt{max\_pos\_weight=30}) against
the single available baseline run sharing the same seed, data split, label
vocabulary, and \emph{current} evaluation schema (\texttt{use\_pos\_weight=
False}). Table~\ref{tab:posweight} reports the matched comparison.

\begin{table}[h]
\centering
\caption{Verified baseline vs.\ positive-class-weighted comparison
(seed 42, matched data/schema).}
\label{tab:posweight}
\begin{tabular}{lccc}
\toprule
Configuration & APS & ROC-AUC & Tuned threshold \\
\midrule
Baseline (no pos.\ weight) & 0.2543 & 0.8911 & 0.15 \\
Positive-class weighted ($\le$30) & 0.2735 & 0.9006 & 0.60 \\
\midrule
Absolute delta & +0.0193 & +0.0096 & -- \\
Relative change & \textbf{+7.6\%} & +1.07\% & -- \\
\bottomrule
\end{tabular}
\end{table}

This verified +7.6\% APS improvement supersedes an earlier, unverified
internal estimate of +4.6\%, which mixed runs from the legacy evaluation
schema. We flag one limitation explicitly: only seed 42 has a full training
run under the current evaluation schema for both the baseline and
weighted arms, so this is a single-seed comparison rather than one averaged
across multiple random seeds; a multi-seed replication is listed as future
work (Section~8).

\subsection{Performance by Clinical Code Type}

Table~\ref{tab:codetype} breaks down test-set performance by the code-type
categories introduced in Section~5.1, weighting each label's metrics by its
test-set support (i.e.\ reflecting performance on an actual occurrence of that
code type, not a per-label average).

\begin{table}[h]
\centering
\caption{Test-support-weighted performance by clinical code type.}
\label{tab:codetype}
\begin{tabular}{lccccc}
\toprule
Code type & \% of test occurrences & Precision & Recall & F1 & AP \\
\midrule
Disease & 85.7\% & 0.222 & 0.298 & 0.200 & 0.220 \\
Administrative & 4.3\% & -- & -- & 0.386 & 0.350 \\
Pregnancy/perinatal & 0.3\% & -- & -- & 0.123 & 0.178 \\
Symptom & 7.1\% & -- & -- & 0.052 & 0.092 \\
Injury & 2.6\% & -- & -- & 0.046 & 0.034 \\
\bottomrule
\end{tabular}
\end{table}

The near-zero performance on injury codes (F1 = 0.046) is consistent with
their nature as unpredictable acute events rather than temporally progressing
conditions, and supports an explicit limitation: the model should not be
expected to anticipate acute trauma or symptom-driven encounters, only
chronic-condition recurrence and, to a lesser extent, administrative/encounter
patterns (which perform surprisingly well, F1 = 0.386, likely reflecting
routine, schedule-driven care patterns).

\subsection{New vs.\ Recurring Diagnosis Performance}

Splitting test predictions by whether the target code is a recurrence (already
in the patient's history) or genuinely new reveals a large asymmetry:
recurring-code predictions achieve F1 = 0.403 (AUC = 0.681), while new-code
predictions achieve F1 = 0.061 (AUC = 0.865). The high AUC but low F1 for new
codes indicates the model \emph{ranks} new diagnoses reasonably well, but the
global decision threshold suppresses nearly all of them, because true new-code
positives are approximately 150$\times$ sparser per sequence position than
recurring-code positives. A threshold sweep confirms that a shared threshold
of 0.50 outperforms the globally tuned 0.60 threshold for \emph{both} subsets
simultaneously, suggesting the global threshold-tuning procedure
over-optimizes for the majority (recurring) subset at the expense of new-code
recall.

\subsection{Per-Support-Tier Threshold Calibration}

Because a single global threshold cannot serve labels with vastly different
base rates well, we evaluated per-support-tier thresholds selected on the
validation split and scored on test. This improves overall F1 from 0.238 to
0.259 and recall from 0.284 to 0.379, and reduces the count of
zero-recall labels from 330 to 305 (out of 469 evaluated labels). Critically,
the ultra-rare support tier remains at F1 = 0 regardless of threshold choice
-- with a mean test support of only $\sim$0.43 occurrences per label, this is
a structural data-scarcity limit that calibration cannot resolve, not a
threshold-tuning problem.

\subsection{Performance by Label Support and Diagnosis-History Length}

Contrary to a naive expectation that performance should degrade monotonically
for rarer labels, per-class ROC-AUC is \emph{not} monotonic with training
support: the ultra-rare support tier achieves a higher mean AUC (0.752) than
the very-common tier (0.673), even though 139 of 469 evaluated labels have
zero recall at the global decision threshold. This apparent paradox is
consistent with AUC's insensitivity to the extreme imbalance within very rare
labels, where a handful of correctly ranked positives can produce a high AUC
despite near-total recall failure at any single operating threshold.

Separately, bucketing test rows by the patient's total historical visit
count (5 tiers, 1--2 through 20+ visits) shows performance
\emph{improving monotonically} with more visit history: sample-wise AUC rises
from 0.892 (1--2 visits) to 0.921 (20+ visits), and F1 rises from 0.223 to
0.256 through the 10--19 tier before a slight dip at 20+. Because the model's
context window is limited to 64 tokens, 37.1\% of test rows exceed this limit
and lose their earliest visits to truncation; counterintuitively, this
truncated subset still \emph{outperforms} the non-truncated subset (F1 = 0.258
vs.\ 0.225, AUC = 0.909 vs.\ 0.896). This suggests that recent history,
which is preserved even under truncation, carries most of the predictive
signal for the (majority) recurrence-driven portion of the task, and that
increasing the context window alone is unlikely to yield large performance
gains without also addressing the structural sparsity of rare labels.

\subsection{Error and Confusion Analysis}

For each label's false positives, we measured the lift of co-occurring true
labels relative to their baseline occurrence rate. Only 17.8\% of the
strongest false-positive confusion pairs fall within the \emph{same} CCSR
category (e.g.\ circulatory-with-circulatory); the large majority are
cross-category and comorbidity-driven (e.g.\ mental/behavioral disorders
co-occurring with pregnancy-related codes, circulatory with neoplasms). This
suggests the model's false positives are substantially explained by genuine
clinical comorbidity patterns -- firing on a label because a clinically
correlated condition is present in the same visit -- rather than superficial
confusion between similarly coded diagnoses.

\subsection{Interpretability Analysis}

\subsubsection{Raw Attention Analysis}

Capturing last-layer, head-averaged \texttt{[CLS]}-token attention over a
500-patient test sample reveals three findings. First, mean attention on
history positions whose code \emph{is} a true next-visit target (0.0251,
$n=7{,}148$) is statistically indistinguishable from attention on
non-target positions (0.0252, $n=9{,}814$) -- the model does not simply
attend more to codes about to recur, a negative result that argues against a
naive ``attention equals importance'' interpretation. Second, attention
exhibits a mild recency bias, peaking at 0.045 for the most recent
pre-\texttt{[CLS]} position and decaying to 0.024 by 19 positions back, even
though Section~5.7 showed that longer total history still improves
performance -- implying that distant-visit information, if used at all, is
aggregated through earlier layers or the residual stream rather than surfaced
directly in final-layer \texttt{[CLS]} attention. Third, \texttt{[SEP]}
(visit-separator) tokens absorb approximately 11\% of total attention mass, a
classic ``attention sink'' pattern that leaves 89\% of attention on
substantive diagnosis codes.

\subsubsection{Integrated Gradients Validation of Confusion Structure}

To test whether the cross-category confusion pattern identified in
Section~5.8 reflects genuine model attribution (rather than only co-occurrence
statistics), we applied Integrated
Gradients~\cite{sundararajan2017axiomatic} to word (diagnosis-code) embeddings
for the 15 labels with the most false positives, attributing each history
code's contribution to the target label's logit. A first pass, ranking
categories by raw attribution mass, produced a misleadingly high 40\%
agreement with the co-occurrence-based confusion categories -- this was found
to be an artifact of category base-rate frequency (common categories such as
circulatory or endocrine dominate every label's naive top-attribution ranking
regardless of any real relationship to that label). Correcting for this by
normalizing attribution mass by each category's baseline occurrence share
(an attribution \emph{lift}, directly analogous to the co-occurrence lift
metric of Section~5.8) yields a corrected agreement rate of \textbf{53.3\%}
(8 of 15 labels), a moderate but real corroboration that the co-occurrence
confusion patterns are partly reflected in actual model attribution rather
than being purely correlational.

\subsubsection{Attention Rollout}

We additionally computed attention
rollout~\cite{abnar2020quantifying}, which recursively combines
head-averaged attention with a residual-aware identity term across all 6
layers, to test whether long-range dependency use is masked by only examining
raw last-layer attention. Contrary to the motivating hypothesis, rollout
attention is \emph{more} concentrated on recent positions than raw last-layer
attention (fraction of mass on the 3 most recent positions: 0.2095 under
rollout vs.\ 0.1932 under raw last-layer attention), the opposite of what
would be expected if distant-history information were being aggregated via
earlier layers and hidden from the final layer. The target-hit-vs-target-miss
null result and the \texttt{[SEP]} attention sink both persist under rollout
(0.0235 vs.\ 0.0250, and 12.0\% \texttt{[SEP]} mass, respectively), confirming
these are robust properties of the model rather than artifacts of examining
only the last layer. The mechanism by which longer diagnosis histories
improve performance (Section~5.7) therefore remains an open question that
neither raw attention nor rollout attention explains.

\section{Discussion}

Taken together, these analyses support a specific, narrow framing of what
this model has learned, and where its limits lie.

\textbf{What the model can credibly claim.} The baseline comparison in
Section~5.2 requires this claim to be narrowed considerably from what the
aggregate headline metrics alone would suggest: a simple logistic-regression
model over flat, non-sequential history features already matches or exceeds
BEHRT on both aggregate ranking metrics (APS, ROC-AUC), and a
gradient-boosted-tree baseline (XGBoost) exceeds it substantially on the same
metrics. BEHRT's evidenced, specific advantage is concentrated in
new-diagnosis detection -- the harder, 30\% minority subset of targets --
where it achieves the highest F1 (0.061) and ROC-AUC (0.865) of any model
tested, including the flat-feature learners, and in producing usable point
predictions (micro-F1 = 0.238) from a single, well-calibrated decision
threshold without requiring per-label tuning, unlike the XGBoost baseline.
Positive-class weighting provides a verified, if single-seed, +7.6\% APS
improvement for rare-code detection on top of this.

\textbf{Where the analogy to general disease prediction breaks.} Prediction
is not the same as mechanistic understanding: the model outputs codes, not an
account of disease pathophysiology, and 70\% of its correct predictions are
recurrences of already-known conditions rather than novel onset detection.
Code-type heterogeneity means 14.3\% of prediction targets are not diseases
at all, and injury/symptom codes -- being driven by acute, largely
unpredictable events -- perform far worse than disease or even administrative
codes. The cohort itself (ICU, multimorbid, median age 62, 0\% pediatric)
limits generalizability to acute or pediatric care settings.

\textbf{Where performance is structurally, not just operationally, limited.}
The ultra-rare label tier remains at F1 = 0 under any threshold, reflecting a
genuine data-scarcity floor (mean test support $<$1 occurrence per label) that
no calibration procedure can overcome; this points toward few-shot or
transfer-learning remedies as the appropriate next step, rather than further
threshold tuning.

\section{Limitations}

\begin{itemize}
  \item \textbf{ICU-specific, adult-only population.} The cohort is not
    representative of general or pediatric populations; external validation
    on other care settings is needed before broader claims can be made.
  \item \textbf{Recurrence-dominated task.} 70\% of targets are recurrences;
    reported aggregate metrics should not be read as a pure measure of
    disease-onset prediction ability.
  \item \textbf{Non-disease target codes.} 14.3\% of targets are symptom,
    administrative, injury, or pregnancy/perinatal codes with substantially
    different (generally lower, except administrative) predictability than
    disease codes.
  \item \textbf{Single-seed pos-weight comparison.} The verified +7.6\% APS
    improvement from positive-class weighting (Section~5.3) is based on one
    matched seed; a multi-seed replication has not yet been performed.
  \item \textbf{Fixed 64-token context window.} 37.1\% of test sequences are
    truncated, though truncation does not currently appear to harm
    performance on the recurrence-dominated portion of the task.
  \item \textbf{Structural rarity floor.} Ultra-rare labels ($<$10 training
    examples) cannot be meaningfully predicted under the current
    architecture and training regime regardless of threshold calibration.
  \item \textbf{Attention is descriptive, not causal.} All attention-based
    analyses (Section 5.9.1, 5.9.3) should be read as exploratory evidence,
    not proof, of the mechanisms underlying model predictions.
  \item \textbf{Untuned baseline calibration.} The XGBoost baseline
    (Section~5.2) uses a single dataset-wide \texttt{scale\_pos\_weight}
    across all 470 labels; we showed that post-hoc per-label threshold
    calibration only marginally improves its micro-F1 (+0.008 absolute),
    so its poor point-prediction F1 despite the best ranking metrics of any
    model tested is more likely driven by the training-time cost-weighting
    itself than by threshold selection -- a fully per-label-weighted refit
    (untested here, see Section~8) remains a lower bound on what a
    fully-tuned flat-feature baseline could achieve, not necessarily its
    ceiling.
\end{itemize}

\section{Future Work}

\begin{enumerate}
  \item \textbf{Multi-seed validation} of the positive-class-weighting effect
    size under the current evaluation schema, to replace the current
    single-seed comparison with a variance-aware estimate.
  \item \textbf{Rare/ultra-rare label remedies}, such as few-shot or
    transfer-learning approaches, since thresholding alone cannot address the
    structural data-scarcity floor identified in Section~5.6.
  \item \textbf{Hierarchical or long-context architectures} to reduce the
    truncation rate for the 37.1\% of patients affected, motivated by the
    observation that longer histories already improve performance even under
    truncation.
  \item \textbf{Recurrence-aware threshold policies} that use separate
    operating thresholds for recurring vs.\ new-code predictions, given the
    threshold-sensitivity asymmetry identified in Section~5.5.
  \item \textbf{Multimodal fusion} of laboratory values, vital signs, and
    clinical notes alongside diagnosis codes.
  \item \textbf{External validation} on non-ICU, non-US, or pediatric
    cohorts to test generalizability of the recurrence-prediction framing.
  \item \textbf{Per-label training-time cost-weighting for flat-feature
    baselines.} We found (Section~5.2) that post-hoc per-label threshold
    calibration only marginally improves XGBoost's point-prediction F1
    (+0.008 absolute), rejecting a purely threshold-based explanation for its
    gap relative to BEHRT. Re-running the XGBoost baseline with per-label
    \texttt{scale\_pos\_weight} set at \emph{training} time (requiring $N$
    separate single-label fits or a custom multi-label objective, rather
    than the one shared weight \texttt{OneVsRestClassifier} currently uses)
    would test whether the remaining gap reflects a training-time
    calibration artifact or a genuine architectural limitation of
    non-sequential models.
\end{enumerate}

\section{Conclusion}

We presented a BEHRT-style transformer for next-visit clinical code
prediction on a 24{,}729-patient, 179{,}713-sample ICU-derived cohort, and --
more importantly -- a systematic, eleven-step audit of what that model is
actually learning. The headline metrics (ROC-AUC 0.901, APS 0.274) are only
meaningful in light of the underlying task composition: 70\% recurrence
versus 30\% onset, 14.3\% non-disease target codes, and a 52{,}559:1 label
imbalance ratio. A baseline comparison against non-learned and flat,
non-sequential learned models further shows that these headline aggregate
metrics can be matched or exceeded by a logistic-regression or XGBoost model
over simple per-patient history features; BEHRT's genuine, defensible
contribution lies specifically in new-diagnosis detection and in yielding a
single well-calibrated decision threshold, not in aggregate ranking
superiority. Positive-class weighting provides a verified, though
single-seed, +7.6\% improvement in rare-code ranking. Two independent
interpretability techniques -- Integrated Gradients and attention rollout --
corroborate that a meaningful share of the model's errors reflect genuine
cross-category clinical comorbidity rather than superficial code confusion,
while also revealing that the mechanism behind long-history benefit remains
unexplained by attention alone. We believe this kind of task-composition
audit, rather than a single AUC/APS pair, is what should determine whether a
BEHRT-style model is framed as a general disease-prediction system or, as we
argue here, a narrower and more defensible clinical decision-support tool for
chronic-condition recurrence and care coordination in multimorbid patients.

\section*{Acknowledgements}

This work was conducted in partnership with the Frost Institute for Data
Science and Computing (IDSC) at the University of Miami. The underlying
critical-care database was used under its applicable data use agreement; the
views expressed are those of the author.

\begin{thebibliography}{99}

\bibitem{li2020behrt}
Y.~Li, S.~Rao, J.~R.~A.~Solares, A.~Hassaine, D.~Canoy, Y.~Zhu, K.~Rahimi, and
G.~Salimi-Khorshidi, ``BEHRT: Transformer for electronic health records,''
\emph{Scientific Reports}, vol.~10, no.~1, pp.~1--12, 2020.

\bibitem{vaswani2017attention}
A.~Vaswani, N.~Shazeer, N.~Parmar, J.~Uszkoreit, L.~Jones, A.~N.~Gomez,
\L.~Kaiser, and I.~Polosukhin, ``Attention is all you need,'' in
\emph{Advances in Neural Information Processing Systems (NeurIPS)}, 2017,
pp.~5998--6008.

\bibitem{devlin2019bert}
J.~Devlin, M.-W.~Chang, K.~Lee, and K.~Toutanova, ``BERT: Pre-training of deep
bidirectional transformers for language understanding,'' in \emph{Proc.
NAACL-HLT}, 2019, pp.~4171--4186.

\bibitem{lee2020biobert}
J.~Lee, W.~Yoon, S.~Kim, D.~Kim, S.~Kim, C.~H.~So, and J.~Kang, ``BioBERT: A
pre-trained biomedical language representation model for biomedical text
mining,'' \emph{Bioinformatics}, vol.~36, no.~4, pp.~1234--1240, 2020.

\bibitem{alsentzer2019clinicalbert}
E.~Alsentzer, J.~R.~Murphy, W.~Boag, W.-H.~Weng, D.~Jin, T.~Naumann, and
M.~B.~A.~McDermott, ``Publicly available clinical BERT embeddings,'' in
\emph{Proc.\ Clinical NLP Workshop, NAACL}, 2019.

\bibitem{rasmy2021medbert}
L.~Rasmy, Y.~Xiang, Z.~Xie, C.~Tao, and D.~Zhi, ``Med-BERT: Pretrained
contextualized embeddings on large-scale structured electronic health records
for disease prediction,'' \emph{npj Digital Medicine}, vol.~4, no.~1,
pp.~1--13, 2021.

\bibitem{choi2016retain}
E.~Choi, M.~T.~Bahadori, J.~A.~Kulas, A.~Schuetz, W.~F.~Stewart, and
J.~Sun, ``RETAIN: An interpretable predictive model for healthcare using
reverse time attention mechanism,'' in \emph{Advances in Neural Information
Processing Systems (NeurIPS)}, 2016.

\bibitem{nguyen2017deepr}
P.~Nguyen, T.~Tran, N.~Wickramasinghe, and S.~Venkatesh, ``Deepr: A
convolutional net for medical records,'' \emph{IEEE Journal of Biomedical and
Health Informatics}, vol.~21, no.~1, pp.~22--30, 2017.

\bibitem{jain2019attention}
S.~Jain and B.~C.~Wallace, ``Attention is not explanation,'' in \emph{Proc.\
NAACL-HLT}, 2019, pp.~3543--3556.

\bibitem{sundararajan2017axiomatic}
M.~Sundararajan, A.~Taly, and Q.~Yan, ``Axiomatic attribution for deep
networks,'' in \emph{Proceedings of the 34th International Conference on
Machine Learning (ICML)}, 2017, pp.~3319--3328.

\bibitem{abnar2020quantifying}
S.~Abnar and W.~Zuidema, ``Quantifying attention flow in transformers,'' in
\emph{Proceedings of the 58th Annual Meeting of the Association for
Computational Linguistics (ACL)}, 2020, pp.~4190--4197.

\bibitem{johnson2023mimiciv}
A.~E.~W.~Johnson, L.~Bulgarelli, L.~Shen, A.~Gayles, A.~Shammout, S.~Horng,
T.~J.~Pollard, S.~Hao, B.~Moody, P.~Raffa, T.~Markazi, R.~G.~Mark, and
R.~Semenov, ``MIMIC-IV, a freely accessible electronic health record
dataset,'' \emph{Scientific Data}, vol.~10, no.~1, p.~1, 2023.

\bibitem{hcup2021ccsr}
Healthcare Cost and Utilization Project (HCUP), ``Clinical Classifications
Software Refined (CCSR) for ICD-10-CM diagnoses,'' Agency for Healthcare
Research and Quality, Rockville, MD, 2021.

\end{thebibliography}

\appendix

\section{Full Training and Run-Control Configuration}

Table~\ref{tab:appendixA} lists the complete run-control configuration for
the reference model (\texttt{clean\_run\_20260910\_124255}), in addition to
the architecture parameters already given in Table~\ref{tab:arch}.

\begin{table}[h]
\centering
\caption{Reference run: full run-control configuration.}
\label{tab:appendixA}
\begin{tabular}{ll}
\toprule
Control & Value \\
\midrule
Epochs & 3 \\
Sample limit (0 = full data) & 0 \\
Gradient-clip norm & 1.0 \\
Early-stop patience & 0 (disabled) \\
Random seed & 42 \\
Use positive-class weighting & True \\
Maximum positive weight & 30.0 \\
Evaluation threshold (default) & 0.5 \\
Threshold tuning enabled & True \\
Threshold-tuning metric & micro F1 \\
Threshold search range & 0.10 -- 0.90 (step 0.05) \\
Selected (tuned) threshold & 0.60 \\
Metrics-excluded labels & \texttt{UNK} \\
\bottomrule
\end{tabular}
\end{table}

\section{Performance by Label Support Tier}

Table~\ref{tab:appendixB} reports mean test-set F1 and ROC-AUC by label
training-support tier, computed on the per-class metrics of the reference
run (469 evaluated labels).

\begin{longtable}{lcccc}
\caption{Performance by training-support tier.}
\label{tab:appendixB} \\
\toprule
Support tier & \# Labels & Mean train support & Mean F1 & Mean ROC-AUC \\
\midrule
\endfirsthead
\toprule
Support tier & \# Labels & Mean train support & Mean F1 & Mean ROC-AUC \\
\midrule
\endhead
Ultra-rare ($<$10) & 38 & $<$10 & $\approx$0.00 & 0.752 \\
Very rare (10--99) & 95 & $\sim$40 & low & -- \\
Rare (100--999) & 92 & $\sim$400 & moderate & -- \\
Common (1{,}000--9{,}999) & 103 & $\sim$3{,}000 & moderate--high & -- \\
Very common ($\ge$10{,}000) & 74 & $\sim$25{,}000 & higher & 0.673 \\
\bottomrule
\end{longtable}

\emph{Note:} exact per-tier mean F1 values for the intermediate tiers are
tracked in the underlying results file
(\texttt{eda/results/08\_disease\_level\_performance\_breakdown.json}); the
key qualitative finding reported in Section~5.7 -- that mean AUC is
\emph{not} monotonic with support, with the ultra-rare tier (0.752)
outperforming the very-common tier (0.673) on AUC despite near-zero F1 -- is
the tier-level result most relevant to this paper's argument and is
reproduced in full precision in Section~5.7.

\section{Performance by Clinical Code Type (Full)}

Table~\ref{tab:appendixC} reproduces the code-type breakdown of
Table~\ref{tab:codetype} together with label-count and training-support
composition, confirming that the target vocabulary's code-type mix (85.6\%
disease / 7.0\% symptom / 4.0\% administrative / 2.7\% injury / 0.7\%
pregnancy, originally estimated from the full unfiltered target distribution)
is closely reproduced when computed directly on the reference run's trained
label vocabulary and test-set occurrences.

\begin{table}[h]
\centering
\caption{Code-type composition and test-support-weighted performance
(reference run, full precision).}
\label{tab:appendixC}
\begin{tabular}{lcccc}
\toprule
Code type & \# Labels & Train support share & Test support share & F1 \\
\midrule
Disease & 326 & 85.55\% & 85.69\% & 0.2004 \\
Symptom & 18 & 7.04\% & 7.06\% & 0.0520 \\
Administrative & 23 & 4.32\% & 4.31\% & 0.3862 \\
Injury & 69 & 2.71\% & 2.61\% & 0.0455 \\
Pregnancy/perinatal & 33 & 0.38\% & 0.34\% & 0.1234 \\
\bottomrule
\end{tabular}
\end{table}

\end{document}
